% Generated by report/make_restated.py from theory/B3_prop59_tightness.tex: do not edit (edit the note and rerun).
\begin{restated}{Remark}{rem:B3-conclusion}{B3}{0}{optimizing $(\gamma,\Delta)$ gains at most a factor $3$}
Within the proof structure of Proposition~5.9 (Lyapunov function $D+\Delta M$, inequalities (I1)--(I5)), optimizing $(\gamma,\Delta)$ changes the certified one-step rate by the factor $Q(s)$, between $1.05$ and $3$, and by at most $4/3$ when $\kappa\ge 9n/16$. This factor is not negligible when $n\gg\kappa$: for $n=1000$ it is $Q(10)\approx2.42$ at $\kappa=10^2$ and $Q(100)\approx2.94$ at $\kappa=10$, of the same size as the speed-ups $3.2$--$3.7$ measured there (Section~\ref{sec:B3-measure}). When $\kappa\le9n/16$, where the textbook rate is $\frac1{3n}$, a factor $3$ is in fact the most that any analysis of $\E[V_t]$, for any $\Delta>0$, can gain over it: by \eqref{eq:B3-EM} and $S_{t-1}\ge0$, $\Ei[M_t]\ge(1-\frac1n)M_{t-1}$, hence $\E[V_t]\ge\Delta(1-\frac1n)^tM_0$, and a bound $\E[V_t]\le c_0\rho^tV_0$ valid for all $t$, with $c_0$ independent of $t$ and $M_0>0$, forces $1-\rho\le\frac1n$. On the hard instances of hypothesis (A4) ($\kappa\gg n$), the rate per effective pass that this proof structure certifies, whatever the constants, is at most $n\tau^\star\le\frac{1.172\,n}{4\kappa}\ll1$. Whenever the exponent of the contraction of $V_t$ observed along SAGA exceeds the certified exponent $-\log\rho_{\mathrm{bound}}$ by more than the factor $Q(s)$, the excess must therefore be located in the state-dependent ratios $r_{\mathrm{sq}},r_{\mathrm{coco}},r_{\mathrm{var}},r_{\mathrm{sc}}$ (inequalities (I1)--(I4)) and in $r_{\min}$, which is exactly what the chains of Proposition~\ref{prop:B3-attribution} measure. The gap itself depends on the instance: on the worst-case instance of caveat (ii) below, for $\kappa\gg n$, it is at most the factor $\frac83$ on the leading term of the complexity, and on the logistic-regression instances of Section~\ref{sec:B3-measure} it comes mainly from (I4) ($\muEff\gg\mu$) and (I2) ($\LEff<R^2$) where it is large ($\kappa\ge10^4$). Two caveats delimit the claim. (i) The optimality of $(\gamma^\star,\Delta^\star)$ is relative to this proof structure: it does not show that $\tau^\star$ is the best rate that an analysis of SAGA can certify under Assumption~\ref{ass:B3-setting}, and another Lyapunov function or another chain of inequalities could certify more. The closed-form analyses we checked give rates of the same order. \cite{hofmann2015variance} give exactly $\tau^\star$ and $\gamma^\star$ (see the paragraph following Proposition~\ref{prop:B3-optimum}); \cite{defazio2014saga} (Theorem~1, Corollary~1 and the ``adaptivity to strong convexity'' comment at the end of the proof of Theorem~1) prove, with a Lyapunov function containing the Bregman divergences $f_i(\phi_i)-f_i(x^\ast)-\inner{f_i'(x^\ast)}{\phi_i-x^\ast}$ and assuming every $f_i$ to be $\mu$-strongly convex, the rate $1-\frac{\mu}{2(\mu n+L)}$ for $\gamma=\frac1{2(\mu n+L)}$ and the rate $1-\min\{\frac1{4n},\frac{\mu}{3L}\}$ for the $\mu$-free step $\gamma=\frac1{3L}$ (this stronger assumption implies Assumption~\ref{ass:B3-setting} with $R^2=L$, and on this subclass the first rate exceeds $\tau^\star$ as soon as $\kappa>n/\sqrt2$, by a factor tending to $2$ as $\kappa/n\to\infty$: an analysis that exploits the stronger assumption does certify more there); when strong convexity only holds on average (our Assumption~\ref{ass:B3-setting}) they state the rate $1-\frac{\mu}{6(\mu n+L)}$ for $\gamma=\frac{1}{3(\mu n+L)}$; the LMI-based analysis of \cite{hu2017unified} (Corollary~3, eq.~(42)) certifies $1-\min\{\frac1{3n},\frac{\mu}{10L}\}$ for $\gamma=\frac1{3L}$ under assumptions implied by Assumption~\ref{ass:B3-setting}: each $f_i$ convex and $L$-smooth ($f_i\in\mathcal F(0,L)$ in their notation) and $F$ in their class $\mathcal S(m,L)$ of functions with sector-bounded gradients, which contains the $m$-strongly convex $L$-smooth functions but, as they stress, also non-convex ones. The rates quoted here that hold under Assumption~\ref{ass:B3-setting} itself are all at most $\tau^\star$, and all of the above have the form $1-\min\{\Theta(\frac1n),\Theta(\frac1\kappa)\}$. (ii) The exponent $1$ of $\kappa$ is intrinsic to SAGA with a constant step, not to its analysis. Proposition~\ref{prop:B3-lower} shows that on an instance of Assumption~\ref{ass:B3-setting} with $n$ identical quadratic components, for every $n$ and every constant step (whether or not it depends on $\mu$), $\E\norm{\theta_t-\etastar}^2$ decays no faster than $(1-\frac{2}{\kappa+1})^{2t}$, so that at least $\frac{\kappa-1}{4}\log\frac1{2\varepsilon}$ iterations are needed, against $O((n+\sqrt{n\kappa})\log\frac1\varepsilon)$ for accelerated methods (hypotheses (A2)--(A4)). For $\kappa\gg n$ the certified complexity $\Theta(\kappa\log\frac1\varepsilon)$ is therefore optimal for SAGA up to a constant factor. The textbook bound guarantees the target after $\frac{16\kappa}{3}\log\frac{1+n/4}{\varepsilon}$ iterations; for the textbook step itself, the identity and the first inequality in \eqref{eq:B3-lower} give on the same instance, with $\theta_0=(r,r)$, $\E\norm{\theta_t-\etastar}^2\ge(1-\frac1{4\kappa})^{2t}\frac12\norm{\theta_0-\etastar}^2$, hence at least $(2\kappa-\frac12)\log\frac1{2\varepsilon}$ iterations (using $\log\frac1{1-x}\le\frac x{1-x}$), which is the bound of part (a) of the SAGA sandwich of the A1 note (Corollary~\ref{cor:saga}) on this instance. On the leading term $\kappa\log\frac1\varepsilon$ the textbook analysis is thus tight up to the factor $\frac83$, and within $\frac{64}{3}\approx21$ of the lower bound valid for every constant step. For oblivious but non-constant step schedules the same conclusion follows from \cite[Theorem~2]{arjevani2016iteration}: ``Suppose that only $L$ the smoothness parameter is known, i.e.\ $\mathcal I(\cdot)=\{L\}$. If the iteration complexity of a given oblivious, possibly stochastic, p-CLI optimization algorithm is $\tilde O(\kappa(f)^\alpha\ln(1/\epsilon))$, then $\alpha\ge1$''; SAGA (with SAG, SVRG and SDCA without duality) is listed there (Section~2.2) as a stationary stochastic p-CLI, and the book's step $\gamma=\frac1{4R^2}$, or any schedule depending on $R^2$ only, is oblivious in their sense (coefficients depending on the side-information $\{L\}$ and on time only). When $\mu$ is also known, their Theorem~1 only gives the lower bound $\tilde\Omega(\sqrt\kappa\ln(1/\epsilon))$ for oblivious schedules, which accelerated gradient descent attains with full gradients, i.e.\ with $n$ component gradients per iteration (for finite sums the oblivious lower bound counted in component gradients is $\tilde\Omega(n+\sqrt{n(\kappa-1)}\ln(1/\epsilon))$ by \cite[Theorem~2]{arjevani2016dimension}, recalled in the A1 note as Theorem~\ref{thm:AS}, and accelerated variance-reduced methods match it for $\kappa\gtrsim n$); this does not contradict Proposition~\ref{prop:B3-lower}, which concerns constant steps. Improving $\kappa$ to $\sqrt{n\kappa}$ therefore requires changing the algorithm (acceleration, hypotheses (A2)--(A3)), consistently with the lower bound of hypothesis (A1). The computer-aided approach of \cite{taylor2019stochastic} automates this kind of search over potential (Lyapunov) functions and parameters by semidefinite programming, with a tightness guarantee for the chosen potential class; it is developed there mainly for sublinear rates of stochastic-gradient-type methods without strong convexity, the authors noting that the same tools apply under strong convexity and that ``minor adaptations'' cover finite-sum methods such as SAGA, and the LMI analysis of \cite{hu2017unified} quoted in (i) is a search of this kind for SAGA with linear rates.
\end{restated}
